\scanpage{127}
\mspage{\handwritten{25}}
\deleted{\handwritten{\deleted{/ obtenues par \unsure{extensions|extension} du corps de base}}

obtient ainsi une \textquotedblleft fibration\textquotedblright\ naturelle de l\unsure{(|'}ouvert dense $U$ des points réguliers de $G$ au dessus du \deleted{\unsure{schéma}} préschéma $\deleted{\unsure{X}}\handwritten{\unsure{\mathscr T|\mathscr C}}$, les fibres \deleted{géométriques} étant \deleted{\illegible{2}} \handwritten{des} ouverts \inserted{denses} de sous-$\deleted{\unsure{S}}$groupes de Cartan des fibres de $G_{\deleted{\unsure{S}}\handwritten{\mathscr T}}$ (savoir les ouverts des points réguliers dans $G$). \deleted{modulo extension du corps de base}. On trouve par exemple\handwritten{\unsure{,|:} le résultat suivant (qui sera considérablement précisé dans l'exposé suivant):}

\uline{Corollaire} 3.3. \quad Soient $G$ un groupe algébrique lisse connexe \deleted{affine} sur le corps $k$, $\handwritten{\mathscr T}$ le schéma des tores maximaux de $G$ ($\handwritten{\simeq}$ le schéma des sous-groupes de Cartan de $G$). Alors le corps des \deleted{\unsure{fnctions|fonctions}} fonctions $k(G)$ \handwritten{de $G$} est isomorphe au corps de fonctions d'un groupe algébrique lisse connexe affine nilpotent $C$ sur le corps des fonctions $k(\handwritten{\mathscr T})$ de $\handwritten{\mathscr T}$, savoir $C=$\deleted{\unsure{1}}\textquotedblleft le sous-groupe de Cartan générique de $G$\textquotedblright. Si \deleted{$G$ est de rang unipotent nul, i.e.} \inserted{\handwritten{$G$ est affine de rang unipotent nul, i.e. si}} \deleted{\unsure{s}}\handwritten{l}es sous-groupes de Cartan \inserted{\handwritten{de $G_{\bar k}$}} sont des tores, alors $k(G)$ est une extension unirationnelle de $k(\handwritten{\mathscr T})$.

Bien entendu, par sous-groupe de Cartan générique de $G$, on entend (par abus de langage) le sous-groupe de Cartan de $G_{k(\handwritten{\mathscr T})}$ fibre générique de $X$ au dessus de $\handwritten{\mathscr T}$. Il n'y a plus qu'à prouver la dernière assertion de 3.3. , qui est contenue dans le \handwritten{résultat bien connu suivant (dû à Chevalley):}

% slanted
\margin{\handwritten{Conserver le texte ci-contre, barré par erreur}}
\deleted{\uline{Lemme} 3.4. \quad Soient $k$ un corps, $T$ un tore sur $k$, $k(T)$ le corps des fonctions rationnelles sur $T$, alors $k(T)$ est une extension unirationnelle de $k$, i.e. est contenue dans une extension transcendante pure de $k$.

Soit en effet $k'$ une extension finie \deleted{\illegible{1}} séparable de $k$ \deleted{\illegible{1}} qui splitte $T$ (Exp X 1.4.), alors \deleted{\unsure{le}} $T\otimes_k k'$ est une variété rationnelle i.e. \deleted{\unsure{isomorphe bir.}} admet un ouvert dense isomorphe à un ouvert dense de l'espace euclidien $\underline E_{\underline k}^{\,n}$, donc $T'=\handwritten{\prod}_{\operatorname{Spec}(k')/\operatorname{Spec}(k)}T_{k'}/\operatorname{Spec}(k')$ est une variété rationnelle (car admet un ouvert dense isomorphe à un ouvert dense de $\handwritten{\prod}_{\operatorname{Spec}(k')/\operatorname{Spec}(k)}E_{k'}^{\,n}$, qui est isomorphe à l'espace \deleted{\unsure{vec}} eucli-}}

% archive sheet number: 126
